\section{Inline Code}
In C the \tscodeinline{strstr} function defined with the prototype
{\ttfamily \PY{k+kt}{char} \PY{o}{*}\PY{n+nf}{strstr}\PY{p}{(}\PY{k}{const} \PY{k+kt}{char} \PY{o}{*}\PY{n}{haystack}\PY{p}{,} \PY{k}{const} \PY{k+kt}{char} \PY{o}{*}\PY{n}{needle}\PY{p}{)}\PY{p}{;}} is
used to locate a substring within a string. It returns a pointer to the first occurrence of the substring
\tscodeinline{needle} in the string \tscodeinline{haystack}, or \tscodeinline{NULL} if the substring is not found.

In Python, you can achieve similar functionality using the \tscodeinline{find} method of strings for example: {\ttfamily \PY{n}{haystack}\PY{o}{.}\PY{n}{find}\PY{p}{(}\PY{n}{sub}\PY{p}{:} \PY{n+nb}{int}\PY{p}{)} \PY{o}{\PYZhy{}\allowbreak{}}\PY{o}{\PYZgt{}} \PY{n+nb}{int}}. This method returns the lowest index of the substring if found in the string, otherwise it returns \tscodeinline{-\allowbreak{}1}.
